\chapter{Introduction}
\label{ch:intro}

\index{acme!introduction}

Acme is a text editor and programming environment created by Rob Pike
at Bell Labs in 1993.  It was designed as part of the Plan~9 operating
system, where it served as both an editor and a shell---a surface on
which all work takes place.

Acme is unlike any other editor you have used.  It has no menus, no
toolbars, no configuration files, no plugin system, and no key bindings
to memorize.  Instead, it offers a blank surface where all commands are
ordinary text that you execute with a mouse click.  Any text, anywhere
on screen, can be a command.  Any text can be a filename.  Any text can
be a search term.  The entire interface is built from three mouse
buttons and text.

\section{Philosophy}

Rob Pike described acme's design philosophy in his 1994 paper:

\begin{quote}
The right model is to have a shell that does not distinguish between a
program that is an editor and a program that is a compiler.  All programs
should have equal access to the services of the environment.
\end{quote}

\index{philosophy}
This leads to several principles:

\begin{description}[style=nextline]
\item[Text is the universal interface.]
Commands are text.  Output is text.  Filenames are text.  You interact
with everything by reading and clicking on text.

\item[The mouse is primary.]
Acme is designed for a three-button mouse.  The keyboard is for typing
text; the mouse is for everything else---executing commands, opening
files, selecting, searching, scrolling, and window management.

\item[Simplicity over features.]
Acme has no syntax highlighting, no auto-completion, no bracket matching,
no built-in debugger, no project manager.  It provides a powerful surface
and lets external programs do the work.

\item[Compose, don't configure.]
Rather than configuring the editor, you compose it with external tools.
Need to sort lines?  Select them and execute \cmd{|sort}.  Need to
reformat?  Execute \cmd{|fmt}.  The Unix toolkit is your extension
mechanism.
\end{description}

\section{How Acme Differs}

If you come from \textbf{vi} or \textbf{Vim}, you will notice that acme
has no modes.  There is no insert mode, no command mode, no visual mode.
You are always editing.  Commands are not key sequences---they are text
you click on.

If you come from \textbf{Emacs}, you will notice that there are no key
bindings to learn beyond basic cursor movement.  There is no Lisp
configuration, no major modes, no minor modes.  The mouse does what
key chords do in Emacs.

If you come from \textbf{VS Code} or similar IDEs, you will notice the
absence of everything: no file tree, no tabs, no status bar, no
integrated terminal widget.  Instead, every window is a general-purpose
text buffer.  A directory listing is a window.  A shell session is a
window.  Compiler output is a window.  They are all the same thing.

The learning curve is not in memorizing commands---there are very few.
The learning curve is in \emph{changing how you think} about interacting
with a computer.

\section{History}

\index{Plan 9}
\index{Pike, Rob}
\index{Bell Labs}
Acme was written by Rob Pike at Bell Labs as part of the Plan~9
operating system.  Plan~9 was the successor to Unix, designed by many
of the same people.  In Plan~9, all resources---files, network
connections, processes, graphics---appear as files in a hierarchical
namespace.  Acme takes full advantage of this: it presents its windows,
buffers, and events as a file system that other programs can read and
write.

\index{plan9port}
\index{Cox, Russ}
In 2003, Russ Cox ported the Plan~9 user-space tools to Unix systems
as \emph{plan9port} (Plan~9 from User Space).  This made acme available
on Linux, macOS, and other Unix systems, though it still depended on
several Plan~9 support programs (devdraw, fontsrv, plumber, 9pserve).

The standalone version documented in this book goes further: it embeds
all those support programs directly into a single binary.  It uses SDL3
for graphics (supporting both X11 and Wayland on Linux, and native
windows on macOS and Windows), FreeType and fontconfig for fonts, and a
built-in plumber for opening files and URLs.  The result is a
single-binary acme with no external dependencies beyond standard system
libraries.

\section{About This Book}

This book is a comprehensive guide to using acme.  It covers
installation on Linux, macOS, and Windows; the mouse and keyboard
interface; all built-in commands; the sam editing language; plumbing;
and practical workflows.

Chapter~\ref{ch:mouse} (The Mouse) is the most important chapter in
this book.  If you read only one chapter, read that one.  Acme is a
mouse-driven editor, and understanding the three buttons and their
chords is essential.

Chapter~\ref{ch:edit} (The Edit Command) covers the most powerful
feature---a complete structural editing language inherited from the sam
editor.  It rewards study but is not required for day-to-day use.

The appendices provide quick-reference tables for commands, mouse
actions, keyboard shortcuts, and the Unicode composition system.

Throughout this book, mouse buttons are referred to as \button{1}
(left), \button{2} (middle), and \button{3} (right).  Keyboard
shortcuts are written as \key{Ctrl+A} (hold Control, press A).
Built-in commands appear as \cmd{Put}, \cmd{Get}, \cmd{New}, etc.
